% Options for packages loaded elsewhere
\PassOptionsToPackage{unicode}{hyperref}
\PassOptionsToPackage{hyphens}{url}
\documentclass[
]{article}
\usepackage{xcolor}
\usepackage[margin=1in]{geometry}
\usepackage{amsmath,amssymb}
\setcounter{secnumdepth}{-\maxdimen} % remove section numbering
\usepackage{iftex}
\ifPDFTeX
  \usepackage[T1]{fontenc}
  \usepackage[utf8]{inputenc}
  \usepackage{textcomp} % provide euro and other symbols
\else % if luatex or xetex
  \usepackage{unicode-math} % this also loads fontspec
  \defaultfontfeatures{Scale=MatchLowercase}
  \defaultfontfeatures[\rmfamily]{Ligatures=TeX,Scale=1}
\fi
\usepackage{lmodern}
\ifPDFTeX\else
  % xetex/luatex font selection
\fi
% Use upquote if available, for straight quotes in verbatim environments
\IfFileExists{upquote.sty}{\usepackage{upquote}}{}
\IfFileExists{microtype.sty}{% use microtype if available
  \usepackage[]{microtype}
  \UseMicrotypeSet[protrusion]{basicmath} % disable protrusion for tt fonts
}{}
\makeatletter
\@ifundefined{KOMAClassName}{% if non-KOMA class
  \IfFileExists{parskip.sty}{%
    \usepackage{parskip}
  }{% else
    \setlength{\parindent}{0pt}
    \setlength{\parskip}{6pt plus 2pt minus 1pt}}
}{% if KOMA class
  \KOMAoptions{parskip=half}}
\makeatother
\usepackage{graphicx}
\makeatletter
\newsavebox\pandoc@box
\newcommand*\pandocbounded[1]{% scales image to fit in text height/width
  \sbox\pandoc@box{#1}%
  \Gscale@div\@tempa{\textheight}{\dimexpr\ht\pandoc@box+\dp\pandoc@box\relax}%
  \Gscale@div\@tempb{\linewidth}{\wd\pandoc@box}%
  \ifdim\@tempb\p@<\@tempa\p@\let\@tempa\@tempb\fi% select the smaller of both
  \ifdim\@tempa\p@<\p@\scalebox{\@tempa}{\usebox\pandoc@box}%
  \else\usebox{\pandoc@box}%
  \fi%
}
% Set default figure placement to htbp
\def\fps@figure{htbp}
\makeatother
\setlength{\emergencystretch}{3em} % prevent overfull lines
\providecommand{\tightlist}{%
  \setlength{\itemsep}{0pt}\setlength{\parskip}{0pt}}
\usepackage{booktabs}
\usepackage{longtable}
\usepackage{array}
\usepackage{multirow}
\usepackage{wrapfig}
\usepackage{float}
\usepackage{colortbl}
\usepackage{pdflscape}
\usepackage{tabu}
\usepackage{threeparttable}
\usepackage{threeparttablex}
\usepackage[normalem]{ulem}
\usepackage{makecell}
\usepackage{xcolor}
\usepackage{siunitx}

    \newcolumntype{d}{S[
      table-align-text-before=false,
      table-align-text-after=false,
      input-symbols={-,\*+()}
    ]}
  
\usepackage{bookmark}
\IfFileExists{xurl.sty}{\usepackage{xurl}}{} % add URL line breaks if available
\urlstyle{same}
\hypersetup{
  hidelinks,
  pdfcreator={LaTeX via pandoc}}

\author{}
\date{\vspace{-2.5em}}

\begin{document}

\section{Introduction}\label{introduction}

L'activité d'assurance repose sur un mécanisme de production inversé :
l'assureur perçoit les primes avant de connaître avec certitude le
montant réel des dépenses qu'il devra engager. En effet, celui-ci est
lié à l'assuré par une promesse d'indemnisation en cas de réalisation du
risque couvert. Le maintien d'un tel système nécessite donc un
renforcement constant des dispositifs de contrôle de la capacité des
assureurs à honorer leurs engagements. C'est dans cette perspective que
s'inscrit le cadre réglementaire de Solvabilité II, entré en vigueur en
2016, dont l'objectif principal est de garantir que les compagnies
d'assurance disposent de fonds propres suffisants pour faire face à
leurs obligations et ainsi instaurer un climat de confiance entre
assureurs et assurés.

La solvabilité d'une compagnie d'assurance désigne sa capacité à
respecter ses engagements financiers à court, moyen et long terme,
notamment en matière de règlement des sinistres et de prestations
contractuelles. Elle constitue ainsi un indicateur fondamental de la
solidité financière des assureurs. Toutefois, cette solvabilité demeure
fortement influencée par l'environnement économique et financier. Par
exemple, une hausse du taux d'inflation peut entraîner une augmentation
mécanique du coût des sinistres, tandis qu'une variation des taux
d'intérêt affecte à la fois les perspectives de rendement des placements
financiers et la valorisation des obligations détenues par les
compagnies d'assurance.

Dans un contexte marqué par la volatilité des marchés financiers, les
risques de perturbation des chaînes de valeur économiques, ainsi que le
renforcement des exigences prudentielles, il devient essentiel de mieux
comprendre les déterminants de la solvabilité des entreprises
d'assurance. Cette problématique est d'autant plus importante que la
stabilité et la pérennité du secteur assurantiel en dépendent
directement.

Le présent travail a donc pour objectif d'évaluer l'impact des variables
financières et macroéconomiques sur la solvabilité des entreprises
d'assurance en France. Plus précisément, il s'agira d'identifier les
variables les plus significatives et de mesurer leur influence sur la
solidité financière des assureurs. La résolution de cette problématique
reposera sur la mobilisation d'outils économétriques appliqués à des
données financières, macroéconomiques et à des indicateurs de
solvabilité propres au secteur de l'assurance.

Cette étude permettra notamment d'apprécier le degré de dépendance de la
solvabilité des assureurs vis-à-vis des variables macro-financières.
Elle constitue ainsi un outil d'aide à la décision pour les compagnies
d'assurance, les décideurs publics ainsi que les autorités de régulation
prudentielle.

\section{Révue de la littérature}\label{ruxe9vue-de-la-littuxe9rature}

Les assureurs jouent un rôle majeur dans le système financier en tant
qu'investisseurs institutionnels et leur exposition aux différents
risques est structurante pour leur solvabilité. Dans le cadre prudentiel
Solvabilité II, la sensibilité des fonds propres aux variations des
taux, des marchés actions et plus généralement de la conjoncture
économique est fortement encadrée. Cette revue de littérature vise à
synthétiser les principaux résultats existants.

\subsubsection{Impact des taux d'intérêt et du marché obligataire sur la
solvabilité des
assureurs}\label{impact-des-taux-dintuxe9ruxeat-et-du-marchuxe9-obligataire-sur-la-solvabilituxe9-des-assureurs}

La longue durée des engagements d'assurance rend ces entreprises
particulièrement sensibles à la structure par terme des taux. Les stress
tests réalisés par l'EIOPA mettent en évidence une forte sensibilité des
entreprises européennes d'assurance aux chocs de marché, révélant
l'exposition directe du SCR aux variations de taux. L'environnement
prolongé de taux bas a exercé une pression sur la solvabilité et la
rentabilité des assureurs notamment à travers leurs fonds propres. La
réévaluation des primes de risques suite à la remontée des taux peut
ainsi entraîner des pertes importantes dans les portefeuilles
obligataires, affectant le ratio de solvabilité (EIOPA, 2018) {[}1{]}.
Dans le cadre de Solvabilité II, la courbe des taux sans risque utilisée
pour actualiser les flux futurs constitue un élément central dans le
calcul des provisions Best Estimates. Un abaissement durable de cette
courbe entraîne une hausse des provisions, particulièrement pour les
passifs de longue duration, ce qui réduit mécaniquement les fonds
propres disponibles (Aalala, 2021) {[}7{]}. Par ailleurs, Berdin et
Gründl (2015) {[}2{]} montrent qu'un environnement de faibles taux peut
conduire les assureurs à accroître leur exposition à des actifs plus
risqués afin de compenser la faiblesse des rendements. Cette
réallocation accroît la sensibilité du bilan aux variations de taux
d'intérêt. Ces travaux confirment que la gestion actif-passif et la
duration des portefeuilles constituent des éléments centraux dans la
maîtrise du risque de solvabilité. Enfin, Poufinas et Siopi (2024)
{[}3{]} mettent en évidence une relation inverse entre le développement
du marché obligataire et le ratio de solvabilité. Un marché obligataire
plus développé se caractérise par une meilleure liquidité et une
meilleure stabilité, conduisant à une hausse de la demande en
obligations et accroit leur prix. Cette appréciation de la valeur des
obligations peut accroître l'exposition au risque de crédit et de
spread, augmentant ainsi le capital réglementaire requis.

\paragraph{Impact du marché actions, de la volatilité et des
corrélations de
marché}\label{impact-du-marchuxe9-actions-de-la-volatilituxe9-et-des-corruxe9lations-de-marchuxe9}

L'exposition des compagnies d'assurance au marché actions s'est accrue
dans un contexte de hausse des taux et d'inflation élevée. Face à cette
remontée des taux, l'EIOPA observe une réallocation partielle des
portefeuilles vers des actifs plus risqués, notamment les actions, afin
de préserver le rendement, au prix d'un besoin accru en capital
réglementaire. Les travaux de Baranoff et Sager (2002) {[}11{]} montrent
qu'une exposition accrue aux risques financiers, notamment actions,
crédit et taux, s'accompagne d'un niveau de capital requis plus élevé
afin d'absorber la volatilité des marchés. Les fluctuations de la valeur
de marché affectent directement les fonds propres et introduisent une
volatilité significative dans le ratio de solvabilité (Rapport ACM 2023)
{[}6{]}. Dans cette logique, les travaux de Pénasse et Dezorme (2021)
{[}5{]} mettent en évidence un arbitrage entre le rendement espéré des
actions et le surplus de SCR imposé par Solvabilité II.

\subsubsection{Rôle des variables macroéconomiques : inflation et cycle
économique}\label{ruxf4le-des-variables-macrouxe9conomiques-inflation-et-cycle-uxe9conomique}

Au-delà des variables financières, la littérature souligne que la
conjoncture macroéconomique, notamment l'inflation et la croissance,
exerce un effet significatif sur la solvabilité des assureurs via la
profitabilité, le niveau des sinistres et la dynamique des primes. Les
travaux récents de l'EIOPA montrent que l'inflation exerce une pression
négative sur la valeur réelle et le rendement des actifs non indexés
ainsi que sur la rentabilité des assureurs. Pour les assureurs non-vie,
elle augmente également le coût des sinistres et des charges
d'exploitation. Pour les assureurs vie, une inflation inattendue peut
rendre les provisions techniques insuffisantes et dégrader la
rentabilité technique, ce qui exerce une pression négative sur la
solvabilité (EIOPA, 2023) {[}9{]}. L'étude de Gonji et al.~(2020)
{[}10{]} met également en évidence un impact du taux d'inflation et du
PIB sur la croissance et la rentabilité des compagnies d'assurance. Une
croissance économique soutenue favorise généralement la demande
d'assurance et la capacité des assureurs à générer des fonds propres,
tandis qu'un ralentissement économique peut dégrader la qualité des
actifs et peser sur le ratio de solvabilité.

Dans l'ensemble, la littérature met en évidence le rôle central des
variables macro-financières dans la dynamique de solvabilité, avec des
effets différenciés selon les branches. Les assureurs Vie apparaissent
généralement plus sensibles aux taux et aux marchés financiers en raison
de la longue duration de leurs engagements, tandis que la branche
Non-Vie semble davantage influencée par l'inflation et le cycle
d'exploitation. Notre travail s'inscrit dans cette lignée tout en s'en
distinguant sur plusieurs points. D'une part, nous mobilisons des
données Solvabilité II françaises sur la période 2016-2025 ce qui permet
de travailler dans un cadre réglementaire homogène peu exploité
empiriquement. D'autre part, nous combinons trois approches
complémentaires --- modèle en différences, ARDL et données de panel ---
afin de capter simultanément les dynamiques de court terme, les
relations d'équilibre de long terme et l'hétérogénéité individuelle des
assureurs. Enfin, la distinction systématique entre branches Vie et
Non-Vie constitue un apport empirique par rapport aux études qui
traitent le secteur de façon agrégée.

\section{Chronologie des
évènements}\label{chronologie-des-uxe9vuxe8nements}

Nous avons réalisé une chronologie des évènements macroéconomiques et
macro-financiers entre 2016 et 2025 (période de nos données) ayant
impacté la solvabilité et l'activité des assureurs. Cela afin de déceler
les variables pertinentes pour notre analyse et afin de déceler des
premiers mécanismes dans l'évolution et la sensibilité du ratio de
solvabilité aux variables macro-financières. Nous avons relevé trois
grandes phases économiques sur cette période. De 2016 à 2020, les
assureurs étaient placés dans un environnement de taux d'intérêt très
bas qui a exercé une pression négative sur leur ratio de solvabilité. En
effet, les revenus obligataires étaient faibles, ce qui a dégradé leur
portefeuille d'investissement et il a été plus difficile pour les
assureurs vie d'honorer leurs engagements. Suite à cela, l'économie a
subi un retournement et un choc inflationniste pendant les années 2021
et 2022 dû à la reprise économique post-COVID-19 et la crise
énergétique. Les taux d'intérêt directeurs de la BCE ont rapidement
augmenté et l'inflation a explosé. Pour les assureurs, la hausse brutale
des taux a entraîné des pertes latentes sur leur portefeuille
obligataire, entraînant une baisse des fonds propres et donc une
dégradation du ratio de solvabilité. Enfin, durant les dernières années
2023 à 2025, l'économie traverse une période d'instabilité financière et
géopolitique. Les marchés financiers sont affectés par plusieurs sources
de volatilité, notamment une crise bancaire et des tensions
géopolitiques toujours dans un contexte inflationniste et de taux
d'intérêt élevés. Les portefeuilles des compagnies d'assurance ont perdu
en valeur et le SCR de marché a augmenté, détériorant les deux
composants du ratio de solvabilité.

Cette chronologie nous a permis de noter les variables macro-économiques
qui semblent en grande partie impacter la solvabilité et l'activité des
assureurs. Ainsi, il a été observé que es taux d'intérêt ont un impact
sur la valeur des obligations détenues par les assureurs entrainant des
pertes ou plus-values latentes réduisant la valeur de marché de leurs
fonds propres. En outre, la rentabilité des portefeuilles obligataires
des assureurs est affectée. Dans ce contexte, les assureurs doivent
ajuster la duration de leur portefeuille afin de gérer leur exposition
au risque de taux. Ces éléments impactent donc le ratio de solvabilité
des assureurs. L'inflation impacte elle aussi significativement
l'activité des assureurs, à travers une hausse du coût des sinistres
pour les assureurs non-vie qui amène les assureurs à ajuster leur
tarification et réévaluer leurs provisions à la hausse. L'inflation
conduit alors une baisse de leur résultat technique et leur solvabilité
est dégradée. Enfin, la volatilité des marchés financiers entraîne une
baisse potentielle de la valeur des portefeuilles actions des assureurs
et une augmentation du SCR de marché pour couvrir cette incertitude.
Cette combinaison réduit les fonds propres et peut dégrader le ratio de
solvabilité.

\section{Formule du calcul du SCR}\label{formule-du-calcul-du-scr}

\section{Synthèse des rapports de conditions de solvabilité et
financière des entreprise
(SFCR)}\label{synthuxe8se-des-rapports-de-conditions-de-solvabilituxe9-et-financiuxe8re-des-entreprise-sfcr}

\section{Statistiques descriptives}\label{statistiques-descriptives}

\subsection{Analyse univariée}\label{analyse-univariuxe9e}

Cette section est consacrée à l'analyse univariée des variables, visant
à décrire leurs distributions et leurs principales caractéristiques
statistiques.

\subsubsection{Répartition des compagnies selon leur
nature}\label{ruxe9partition-des-compagnies-selon-leur-nature}

L'activité d'assurance non-vie concentre la majorité des compagnies du
marché. À l'inverse, les Fonds de Retraite et de Prévoyance Sociale
(FRPS) demeurent très faiblement représentés, avec une part limitée à
\(0,91 \%\). De même, le segment de la réassurance ne regroupe qu'une
fraction marginale des compagnies d'assurance, soit \(2,56 \%\).

\begin{center}\includegraphics[width=0.6\linewidth]{Stat_desc_files/figure-latex/unnamed-chunk-2-1} \end{center}

\subsubsection{Répartition des compagnies par
famille}\label{ruxe9partition-des-compagnies-par-famille}

La structure du marché est majoritairement dominée par les organismes
mutualistes, lesquels représentent \(53,71 \%\) de l'ensemble des
compagnies. Le solde du marché se répartit de manière relativement
équilibrée entre les assureurs traditionnels et les bancassureurs.

\begin{center}\includegraphics[width=0.6\linewidth]{Stat_desc_files/figure-latex/unnamed-chunk-3-1} \end{center}

\subsection{Analyse bivarariée}\label{analyse-bivarariuxe9e}

À l'exception du segment de la réassurance, les organismes mutualistes
apparaissent comme majoritairement représentés dans l'ensemble des
regroupements par nature d'activité. Leur part relative demeure plus
élevée que celle des bancassureurs et des assureurs traditionnels, tant
dans les catégories des prestataires d'assurance non-vie que dans celles
de l'assurance vie et de la capitalisation. En revanche, le secteur de
la réassurance est dominé par les assureurs traditionnels, lesquels
représentent \(71,7\%\) des acteurs recensés. Il convient également de
souligner que seuls les organismes mutualistes et les assureurs
traditionnels interviennent dans la fourniture de fonds de retraite et
de prévoyance sociale (FRSP). Enfin, la répartition des différentes
familles d'assureurs apparaît globalement équilibrée dans les activités
mixtes combinant assurance vie et assurance non-vie.

Dans un soucis d'avoir des données de taille comparable entre les
catégorie mixte, non-vie et vie, nous avons fusionné les catégories
mixte et vie.

\begin{center}\includegraphics[width=0.6\linewidth]{Stat_desc_files/figure-latex/unnamed-chunk-4-1} \end{center}

\subsection{Analyses trimestrielles}\label{analyses-trimestrielles}

Pour toutes les analyses qui suivent, nous ne prendrons en compte que
deux natures de sociétés:

\begin{itemize}
\tightlist
\item
  Non Vie;
\item
  Vie (regroupant les natures \textbf{Vie/Capitalisation} et
  \textbf{Mixte}).
\end{itemize}

La base de données contient 3444 trimestres avec des variables cibles
déclarées. Nous pouvons donc observer l'évolution moyenne des Fonds
propres et du ratio de solvabilité avec les données disponibles.

\begin{enumerate}
\def\labelenumi{(\roman{enumi})}
\tightlist
\item
  \textbf{Ratio de solvabilité}
\end{enumerate}

\begin{center}\includegraphics[width=0.6\linewidth]{Stat_desc_files/figure-latex/unnamed-chunk-5-1} \end{center}

Le ratio moyen de solvabilité trimestriel apporte plus d'information sur
les variations qu'ont connues les sociétés d'assurance au cours de
l'année. On constate donc que le ratio de solvabilité moyen des sociétés
oscille avec beaucoup plus d'amplitude que les données annuelles ne le
laissent entrevoir. Le ratio moyen trimestriel des sociétés Non Vie
varie avec une tendance légèrement baissière entre \textbf{320\(\%\)} et
\textbf{250\(\%\) }. Il y a également une volatilité qui est
décroissante dans le temps. Quant au ratio trimestriel moyen des
sociétés d'assurances Vie, il est globalement plus faible que celui des
sociétés non vie évoluant entre \textbf{160\(\%\)} et
\textbf{260\(\%\)}. Entre 2018 et 2020, on remarque une
désynchronisation très forte entre le ratio moyen annuel et trimestriel
des sociétés d'assurances vie.

\begin{enumerate}
\def\labelenumi{(\roman{enumi})}
\setcounter{enumi}{1}
\tightlist
\item
  \textbf{Fonds Propres}
\end{enumerate}

\begin{center}\includegraphics[width=0.6\linewidth]{Stat_desc_files/figure-latex/4-1} \end{center}

Les données trimestrielles offrent une vision plus détaillée de
l'évolution des fonds propres. Ces graphiques montrent que les fonds
propres des sociétés d'assurance Vie sont plus volatils que ceux des
sociétés Non‑Vie. Par ailleurs, les fonds propres moyens des sociétés
Vie et Non‑Vie semblent présenter une légère tendance haussière, ce qui
suggère que les entreprises de notre échantillon accumulent
progressivement davantage de fonds propres au fil du temps, malgré
certaines fluctuations.

Nous observons également un écart important entre les valeurs annuelles
et trimestrielles, aussi bien pour les sociétés d'assurance Non‑Vie que
pour les sociétés Vie. Ce biais pourrait provenir de la taille des
entreprises déclarant des données trimestrielles. En effet, toutes les
sociétés ne transmettent pas d'informations trimestrielles. Seules 131
sociétés de la base disposent d'au moins une observation trimestrielle,
ce qui représente environ 50.38\(\%\) de l'ensemble des sociétés
recensées.

\begin{center}\includegraphics[width=0.6\linewidth]{Stat_desc_files/figure-latex/Décalage-1} \end{center}

La divergence entre les moyennes annuelles et trimestrielles peut
s'expliquer par le fait que les sociétés ne déclarant pas de données
trimestrielles présentent, en moyenne, des niveaux de fonds propres
annuels plus faibles. Autrement dit, seules les grandes sociétés
d'assurance transmettent des données trimestrielles, tandis que toutes
les sociétés --- y compris les plus petites --- déclarent des données
annuelles. Ces petites sociétés tirent donc la moyenne annuelle globale
vers le bas, ce qui crée un écart entre les deux séries.

Pour corriger ce biais de sélection présent dans les données
trimestrielles, nous allons ajuster les séries trimestrielles à partir
des données annuelles.

\paragraph{Ajustement des valeurs
trimestrielles}\label{ajustement-des-valeurs-trimestrielles}

L'objectif est d'estimer le montant total de SCR et de fonds propres
trimestriels que nous aurions obtenu si toutes les sociétés avaient
transmis des valeurs trimestrielles. Pour cela, pour les sociétés ne
déclarant que des valeurs annuelles, nous estimerons leurs valeurs
trimestrielles en nous appuyant sur les comportements observés chez les
sociétés qui déclarent à la fois des données trimestrielles et
annuelles.

La méthode d'ajustement repose sur le calcul du taux de croissance moyen
par nature, obtenu à partir des sociétés qui déclarent à la fois des
valeurs annuelles et trimestrielles. Ces taux de croissance moyens
servent ensuite à estimer les valeurs trimestrielles manquantes à partir
des valeurs annuelles disponibles.

\begin{enumerate}
\def\labelenumi{(\roman{enumi})}
\tightlist
\item
  Pour le premier trimestre
\end{enumerate}

Le taux de croissance moyen utilisé pour l'ajsutement correspond au taux
de croissance entre la somme des valeurs de l'année précédente et la
somme des valeurs du premier trimestre de l'année suivante.

\begin{description}
\tightlist
\item[Supposons que nous souhaitons déterminer le taux de croissance
moyen du premier trimestre de l'année t. Pour cela, nous considèrons
l'ensemble E des sociétés ayant déclaré les valeurs de l'année t-1 ainsi
que celles du premier trimestre de l'année t. Le taux de croissanse se
calcul donc de la manière suivante :]
\end{description}

\[ \tau_{moyen_{1,t}} =   \frac{\sum_{j \in E} SCR^j_{T_{1,t}}-\sum_{j \in E} SCR^j_{A_{t-1}} }{\sum_{j \in E}SCR^j_{A_{t-1}}} \]
où :

\begin{itemize}
\tightlist
\item
  \(SCR^j_{T_{1,t}}\) désigne le SCR de la société j au premier
  trimestre de l'année t;
\item
  \(SCR^j_{A_{t-1}}\) désigne le SCR de la société j de l'année t-1.
\end{itemize}

Nous faisons l'hypothèse que les sociétés n'ayant pas déclaré le premier
trimestre de l'année t auraient connu un taux de croissance
\(\tau _{moyen_{1,t}}\) de leur SCR entre la fin de l'année t‑1 et le
premier trimestre de l'année t. Nous estimons donc le SCR des sociétés
n'ayant pas déclaré au premier trimestre de l'année t en multipliant la
somme de leur SCR de l'année t‑1 par \(1+\tau _{moyen_{1,t}}\). Pour
estimer le SCR moyen au premier trimestre de l'année t, il suffit alors
de considérer la somme des SCR des sociétés ayant déclaré, d'y ajouter
la somme estimée pour celles qui n'ont pas déclaré, puis de diviser
l'ensemble par le nombre total de sociétés.

\begin{enumerate}
\def\labelenumi{(\roman{enumi})}
\setcounter{enumi}{1}
\tightlist
\item
  Pour les autres trimestres (T2, T3 et T4)
\end{enumerate}

Le taux de croissance moyen utilisé pour l'imputation correspond au taux
de croissance entre la somme des valeurs du trimestre précédent et la
somme des valeurs du trimestre que l'on souhaite estimer. Les valeurs
estimer du premier trimestre servent à estimer ceux du deuxième, les
valeurs de deuxième servent à estimer le troisième et les valeurs du
troisième servent à estimer ceux du quatrième.

Supposons que nous souhaitons déterminer le taux de croissance moyen du
trimestre i\(\in \{2,3,4\}\) de l'année t. Pour cela, nous considèrons
l'ensemble \(E'\) des sociétés ayant déclaré les valeurs aux trimestres
i-1 et i de l'année t. Le taux de croissanse se calcul donc de la
manière suivante :

\[ \tau_{moyen_{i,t}} =   \frac{\sum_{j \in E'} SCR^j_{T_{i,t}}-\sum_{j \in E'} SCR^j_{T_{i-1,t}} }{\sum_{j \in E'}SCR^j_{T_{i-1,t}}} \]
Nous estimons le SCR des sociétés n'ayant pas déclaré au trimestre i de
l'année t en multipliant l'estimation de leur SCR du trimestre i-1 de
l'année t par \(1+\tau _{moyen_{i,t}}\). Pour estimer le SCR moyen au
trimestre i de l'année t, il suffit alors de considérer la somme des SCR
des sociétés ayant déclaré, d'y ajouter la somme estimée pour celles qui
n'ont pas déclaré, puis de diviser l'ensemble par le nombre total de
sociétés.

\subsubsection{Visualisation de
l'ajustement}\label{visualisation-de-lajustement}

\begin{center}\includegraphics[width=0.6\linewidth]{Stat_desc_files/figure-latex/8-1} \end{center}

Après ajustement, on observe que les données trimestrielles s'alignent
désormais mieux avec les données annuelles, ce qui confirme la
correction du biais initialement identifié.

\section{Modélisation}\label{moduxe9lisation}

La complexité des interactions entre les variables macro-financières et
la solvabilité des assureurs nécessite une approche empirique
séquentielle. Notre stratégie de modélisation s'articule autour de trois
axes complémentaires :

\begin{enumerate}
  \item \textbf{}Modèle en données de panel : Cette première modélisation vise à analyser la dynamique globale des indicateurs de solvabilité tout en tenant compte de l'hétérogénéité individuelle des assureurs.
  \item \textbf{}Régressions en différence simple : Nous effectuons une régression MCO des différences logarithmiques en ne retenant que la dimension temporelle afin d'évaluer la dynamique de court terme.
  \item \textbf{}Modèle ARDL (Auto Regressive Distributed Lag) : Il permet d'intégrer simultanément des relations de court terme et d'éventuelles relations de long terme (cointégration) qui ne peuvent pas être captées par les modèles de régression simple.
\end{enumerate}

\subsection{I - Modèle de données de
panel}\label{i---moduxe8le-de-donnuxe9es-de-panel}

\subsection{II - Régressions en
différence}\label{ii---ruxe9gressions-en-diffuxe9rence}

Les régressions en différence constituent étape d'analyse des liens de
court terme entre les variables macro-financières et les indicateurs de
solvabilité. Cette approche permet d'étudier la transmission immédiate
des chocs macro-financiers tout en intégrant une dynamique
autorégressive et d'éventuels effets retardés. Nous utilisons des
modèles de régression linéaire estimés par les Moindres Carrés
Ordinaires. La spécification générale du modèle est la suivante :

\begin{equation}
\Delta Y_t = \alpha + \sum_{i=1}^{n} \beta_i \Delta X_{i, t-k} + \rho \Delta Y_{t-1} + \epsilon_t
\end{equation}

Où :

\begin{itemize}
    \item $\Delta Y_t$ représente le taux de croissance de l'indicateur de solvabilité étudié (Ratio, Fonds Propres ou SCR) au temps $t$.
    \item $\alpha$ est la constante du modèle (ou drift).
    \item $\Delta X_{i, t-k}$ représente la variation de la $i$-ème variable explicative (Indices boursiers, IPC, Taux, PIB) avec un retard $k \geq 0$.
    \item $\rho$ est le coefficient associé au terme autorégressif, captant l'inertie de la variable expliquée.
    \item $\epsilon_t$ est le terme d'erreur (bruit blanc) $\epsilon_t \sim N(0, \sigma^2)$.
\end{itemize}

Ce modèle repose sur les hypothèses classiques des MCO : linéarité de la
relation, exogénéité des régresseurs, homoscédasticité et absence
d'autocorrélation des résidus, et normalité de la distribution des
erreurs. Ces hypothèses sont vérifiées empiriquement en fin de partie à
l'aide des tests appropriés.

\subsubsection{1) Agrégation et transformation des
données}\label{agruxe9gation-et-transformation-des-donnuxe9es}

Dans le cadre de cette modélisation, les données individuelles des
assureurs et les variables macro-financières ont été agrégées par
trimestre afin de ne conserver qu'une dimension temporelle. Cette
agrégation permet d'obtenir une série synthétique reflétant l'évolution
globale du secteur et de s'aligner sur les dates de reporting
réglementaires du secteur. Ce choix méthodologique implique dans un
premier temps de renoncer à l'analyse des effets individuels mais il est
cohérent avec l'objectif de l'étude, centré sur la dynamique
macroéconomique de la solvabilité.

Nous modéliserons les fonds propres, le SCR et le ratio de solvabilité
des assureurs. Nous avons retenu les variables explicatives suivantes
(valeurs trimestrielles) : les indices boursiers S\&P 500 et Euro Stoxx
50, les taux d'intérêt de court et long terme, l'indice des prix à la
consommation en France (IPC) et le Produit Intérieur Brut français.

Pour chaque indice boursier (S\&P 500, EURO STOXX 50, DAX, CAC 40), nous
effectuons deux transformations majeures pour assurer la pertinence de
l'analyse :

\begin{itemize}
  \item \textbf{Agrégation trimestrielle} : Afin de faire correspondre les données boursières aux dates de remise des rapports de solvabilité, nous retenons la dernière valeur de clôture enregistrée pour chaque trimestre.
  \item \textbf{Calcul des rendements logarithmiques} : Nous utilisons la formule $r_t = \ln(P_t) - \ln(P_{t-1})$. Les log-rendements sont préférés car ils permettent de traiter des séries temporelles stationnaires et facilitent la comparaison entre indices de niveaux de prix différents.
\end{itemize}

\subsubsection{2) Spécification des variables et
stationnarité}\label{spuxe9cification-des-variables-et-stationnarituxe9}

Nous avons également enrichi nos données en calculant les différences
logarithmiques des variables qui seront utilisées dans la régression en
différence. Cette transformation permet de rendre les séries
stationnaires et d'interpréter les variations comme des taux de
croissance. Les taux d'intérêt sont quant à eux exprimés en différences
premières, correspondant à des variations en niveau.

L'ensemble des variables sont transformées en différences logarithmiques
selon \(\Delta\ln(X_t) = \ln(X_t) - \ln(X_{t-1})\), à l'exception des
taux d'intérêt exprimés en différences premières
\(\Delta X_t = X_t - X_{t-1}\), ces derniers pouvant prendre des valeurs
négatives incompatibles avec le logarithme.

Afin de pouvoir mener des régressions robustes, il est nécessaire que
l'ensemble de nos séries utilisées soient stationnaires. La
stationnarité des séries transformées a été vérifiée à l'aide des tests
de Dickey-Fuller augmenté (ADF) et KPSS. Le test ADF repose ici sur
l'observation de la p-value associée au test :

\begin{itemize}
    \item $H_0$ : La série possède une racine unitaire (non-stationnaire)
    \item $H_1$  : La série est stationnaire
\end{itemize}

Si \(t_{stat}\) \textless{} \(t_{crit}\) : On rejette \emph{H0} au seuil
de significativité choisi.

Les hypothèses du test KPSS sont les suivantes :

\begin{itemize}
  \item $H_0$ : La série est stationnaire
  \item $H_1$ : La série est non stationnaire
\end{itemize}

Si \(t_{stat}\) \textgreater{} \(t_{crit}\) : On rejette \emph{H0} au
seuil de significativité choisi.

Ainsi, les tests ADF et KPSS appliqués aux séries transformées
confirment leur stationnarité validant les conditions d'application du
modèle de régression linéaire Moindres Carrés Ordinaires.

\subsubsection{3) Test de Causalité de
Granger}\label{test-de-causalituxe9-de-granger}

Des tests de causalité au sens de Granger ont été réalisés afin
d'identifier les retards pertinents des variables macro-financières dans
les modèles dynamiques. Les tests ont été conduits sur des horizons
allant de un à quatre trimestres afin de capturer d'éventuels délais de
transmission. L'hypothèse nulle du test est la suivante :

\emph{H0} : les coefficients associés aux valeurs retardées de la
variable explicative sont conjointement nuls.

Pour les assureurs Vie, les résultats mettent en évidence un pouvoir
prédictif significatif des taux longs (retard 1), de l'IPC (retard 4) et
du S\&P500 (retards 2) sur les fonds propres. L'IPC (retard 4) a un
effet prédictif significatif sur le SCR et le S\&P500 (retard 3) sur le
ratio. Pour les assureurs Non-Vie, seules quelques relations prédictives
limitées apparaissent, principalement entre l'IPC (retard 3) et les
fonds propres ainsi qu'entre les taux longs (retard 1) et le ratio de
solvabilité.

\subsubsection{4) Estimation du modèle}\label{estimation-du-moduxe8le}

Afin d'éviter les problèmes de multicolinéarité entre indices boursiers
fortement corrélés, nous avons retenu uniquement le S\&P 500 et l'Euro
Stoxx. Ce choix permet de conserver une représentation pertinente des
dynamiques des marchés, en combinant un indice global et un indice
européen ce qui est cohérent avec l'exposition des assureurs étudiés.

\paragraph{\texorpdfstring{4.1) Spécifications\\
}{4.1) Spécifications }}\label{spuxe9cifications}

\hfill\break
Des spécifications contemporaines ont d'abord été estimées par MCO
(cf.~Annexe Table). Malgré des diagnostics résiduels satisfaisants, les
coefficients apparaissaient majoritairement non significatifs, suggérant
l'existence de mécanismes de transmission retardés. Les modèles
dynamiques finaux ont donc été construits en deux étapes : les tests de
causalité de Granger ont d'abord permis d'identifier les retards
pertinents puis les spécifications ont été comparées via les critères
AIC et BIC afin d'assurer un équilibre entre robustesse statistique et
interprétation économique.

Les résultats mettent en évidence des mécanismes de transmission
différenciés entre branches Vie et Non-Vie ce qui est cohérent avec
leurs structures d'activité respectives.

Les écarts-types des coefficients sont corrigés à l'aide de l'estimateur
robuste de Newey-West afin de tenir compte d'une éventuelle
hétéroscédasticité et autocorrélation des résidus. Bien que les tests de
Breusch-Pagan et Ljung-Box ne rejettent pas formellement
l'homoscédasticité et l'absence d'autocorrélation, la sensibilité des
p-values à cette correction suggère la présence de déviations non
détectées par ces tests sur petit échantillon, justifiant une approche
robuste par précaution.

\subparagraph{4.2.1) Assureurs Vie}\label{assureurs-vie}

\subparagraph{Fonds propres - Vie}\label{fonds-propres---vie}

\begin{equation}
\begin{aligned}
\Delta \ln(\text{FP}_t) = \alpha & + \rho \Delta \ln(\text{FP}_{t-1})
& + \beta_1 \Delta \text{TauxL}_{t-1}
& + \beta_2 \Delta \ln(\text{IPC})_{t-4}
& + \beta_3 \Delta \ln(\text{S\&P500})_{t-2} + \epsilon_t
\end{aligned}
\end{equation}

\begin{table}[H]
\centering\centering
\caption{\label{tab:unnamed-chunk-22}Analyse dynamique des Fonds Propres - Branche Vie}
\centering
\begin{tabular}[t]{lccc}
\toprule
\multicolumn{1}{c}{ } & \multicolumn{3}{c}{Fonds Propres (Vie)} \\
\cmidrule(l{3pt}r{3pt}){2-4}
  & Est. & t & p\\
\midrule
\cellcolor{gray!10}{(Intercept)} & \cellcolor{gray!10}{\num{0.0352}} & \cellcolor{gray!10}{\num{1.4257}} & \cellcolor{gray!10}{\num{0.1650}}\\
dplyr = lag(diff\_fonds\_propres, 1) & \num{-0.5197}* & \num{-2.6712} & \num{0.0124}\\
\cellcolor{gray!10}{dplyr = lag(diff\_taux\_long, 1)} & \cellcolor{gray!10}{\num{-0.0901}**} & \cellcolor{gray!10}{\num{-3.1146}} & \cellcolor{gray!10}{\num{0.0042}}\\
dplyr = lag(rendement\_sp500, 2) & \num{0.3673}** & \num{3.0802} & \num{0.0046}\\
\cellcolor{gray!10}{dplyr = lag(diff\_ipc, 4)} & \cellcolor{gray!10}{\num{-4.4542}*} & \cellcolor{gray!10}{\num{-2.5296}} & \cellcolor{gray!10}{\num{0.0173}}\\
\midrule
Num.Obs. & \num{33} &  & \\
\cellcolor{gray!10}{R2} & \cellcolor{gray!10}{\num{0.390}} & \cellcolor{gray!10}{} & \cellcolor{gray!10}{}\\
R2 Adj. & \num{0.303} &  & \\
\bottomrule
\end{tabular}
\end{table}

L'ensemble des coefficients sont significatifs au seuil de 5\% après
correction des écarts-types par la méthode de Newey-West. Les résidus ne
présentent pas d'autocorrélation selon le test de Ljung-Box, le modèle
semble donc robuste. Les résultats mettent en évidence une forte
sensibilité des fonds propres des assureurs Vie aux conditions
macro-financières. La hausse des taux d'intérêt de long terme exerce un
effet négatif sur la croissance des fonds propres traduisant
probablement la baisse de la valeur de marché des portefeuilles
obligataires lors des phases de remontée des taux. À l'inverse, les
performances des marchés actions internationaux, approximées par le S\&P
500, soutiennent positivement l'évolution des fonds propres avec un
décalage de deux trimestres confirmant l'exposition importante des
assureurs Vie aux marchés financiers. L'inflation apparaît également
comme un facteur défavorable à la solvabilité avec un effet retardé de
quatre trimestres. Ce résultat peut refléter l'augmentation progressive
des coûts de gestion et le délai d'ajustement des primes. Enfin, la
significativité du terme autorégressif souligne l'existence d'une
dynamique de persistance des fonds propres, les chocs observés tendent à
se diffuser progressivement dans le temps.

\subparagraph{SCR - Vie}\label{scr---vie}

\begin{equation}
\begin{aligned}
\Delta \ln(\text{SCR}_t) = \alpha & + \rho \Delta \ln(\text{SCR}_{t-1}) 
& + \beta_1 \Delta \text{TauxL}_{t-1} 
& + \beta_2 \Delta \ln(\text{IPC})_{t-3} + \epsilon_t
\end{aligned}
\end{equation}

\begin{table}[H]
\centering\centering
\caption{\label{tab:unnamed-chunk-23}Analyse dynamique du SCR - Branche Vie}
\centering
\begin{tabular}[t]{lccc}
\toprule
\multicolumn{1}{c}{ } & \multicolumn{3}{c}{SCR (Vie)} \\
\cmidrule(l{3pt}r{3pt}){2-4}
  & Est. & t & p\\
\midrule
\cellcolor{gray!10}{(Intercept)} & \cellcolor{gray!10}{\num{0.0388}**} & \cellcolor{gray!10}{\num{2.9627}} & \cellcolor{gray!10}{\num{0.0058}}\\
dplyr = lag(diff\_scr, 1) & \num{-0.4102}** & \num{-2.7535} & \num{0.0098}\\
\cellcolor{gray!10}{dplyr = lag(diff\_taux\_long, 1)} & \cellcolor{gray!10}{\num{-0.0873}**} & \cellcolor{gray!10}{\num{-2.9415}} & \cellcolor{gray!10}{\num{0.0061}}\\
dplyr = lag(diff\_ipc, 3) & \num{-4.0140}** & \num{-3.3197} & \num{0.0023}\\
\midrule
\cellcolor{gray!10}{Num.Obs.} & \cellcolor{gray!10}{\num{35}} & \cellcolor{gray!10}{} & \cellcolor{gray!10}{}\\
R2 & \num{0.352} &  & \\
\cellcolor{gray!10}{R2 Adj.} & \cellcolor{gray!10}{\num{0.289}} & \cellcolor{gray!10}{} & \cellcolor{gray!10}{}\\
\bottomrule
\end{tabular}
\end{table}

Les résultats montrent que le SCR des assureurs Vie présente une
dynamique de rappel importante comme l'indique la significativité du
terme autorégressif. Les ajustements du besoin en capitals sont donc
progressifs et les chocs observés sur le SCR tendant à se résorber
partiellement au trimestre suivant. Les taux d'intérêt de long terme et
l'inflation exercent tous deux un effet négatif sur la croissance du
SCR. Ces résultats apparaissent cohérents dans la mesure où ces deux
variables influencent les mécanismes d'actualisation des engagements des
assureurs. Dans un contexte de hausse des taux, souvent associé à des
tensions inflationnistes, la valeur actuelle des passifs futurs diminue,
ce qui contribue à alléger les exigences réglementaires en capital.
L'effet retardé de l'IPC suggère toutefois que cette transmission
s'opère progressivement dans le bilan des assureurs Vie. Par ailleurs,
la constante positive et significative indique l'existence d'une
croissance résiduelle du SCR non entièrement expliquée par les variables
macro-financières intégrées au modèle.

\subparagraph{Ratio de solvabilité -
Vie}\label{ratio-de-solvabilituxe9---vie}

\begin{equation}
\begin{aligned}
\Delta \ln(\text{Ratio\_Moyen}_t) = \alpha & + \rho \Delta \ln(\text{Ratio\_Moyen}_{t-1})
& + \beta_1 \Delta \ln(\text{S\&P500})_{t-3} + \epsilon_t
\end{aligned}
\end{equation}

\begin{table}[H]
\centering\centering
\caption{\label{tab:unnamed-chunk-24}Analyse dynamique du ratio de solvabilité - Branche Vie}
\centering
\begin{tabular}[t]{lccc}
\toprule
\multicolumn{1}{c}{ } & \multicolumn{3}{c}{Ratio (Vie)} \\
\cmidrule(l{3pt}r{3pt}){2-4}
  & Est. & t & p\\
\midrule
\cellcolor{gray!10}{(Intercept)} & \cellcolor{gray!10}{\num{-0.0130}} & \cellcolor{gray!10}{\num{-1.3896}} & \cellcolor{gray!10}{\num{0.1742}}\\
dplyr = lag(diff\_ratio, 1) & \num{-0.4683}*** & \num{-4.3912} & \num{0.0001}\\
\cellcolor{gray!10}{dplyr = lag(rendement\_sp500, 3)} & \cellcolor{gray!10}{\num{0.2843}+} & \cellcolor{gray!10}{\num{1.7086}} & \cellcolor{gray!10}{\num{0.0972}}\\
\midrule
Num.Obs. & \num{35} &  & \\
\cellcolor{gray!10}{R2} & \cellcolor{gray!10}{\num{0.307}} & \cellcolor{gray!10}{} & \cellcolor{gray!10}{}\\
R2 Adj. & \num{0.263} &  & \\
\bottomrule
\end{tabular}
\end{table}

La dynamique du ratio de solvabilité semble principalement expliquée par
son propre terme retardé dont le coefficient négatif traduit un
mécanisme de retour à la moyenne dans l'ajustement de la solvabilité des
assureurs Vie. Le rendement du S\&P500 exerce un effet positif modéré
sur le ratio de solvabilité avec un retard de trois trimestres,
suggérant qu'un environnement boursier favorable améliore
progressivement la solvabilité des assureurs Vie à travers la
valorisation des actifs financiers et de meilleures conditions de
rendement des placements. L'absence d'effet significatif sur le SCR
suggère que cette amélioration des fonds propres ne semble pas être
entièrement compensée par une hausse des exigences en capital. À
l'inverse, les variables de taux d'intérêt et d'inflation n'apparaissent
plus comme des déterminants robustes des variations du ratio de
solvabilité probablement en raison d'effets compensatoires entre fonds
propres et SCR. Par ailleurs, la constante n'apparaît plus
significative, il ne semble pas y avoir de croissance structurelle du
ratio non expliquée par les variables macro-financières et la dynamique
de rappel.

\paragraph{\texorpdfstring{4.2.2) Assureurs Non-Vie\\
}{4.2.2) Assureurs Non-Vie }}\label{assureurs-non-vie}

\hfill\break
Nous adoptons la même démarche pour les assureurs Non-Vie malgré
l'absence de résultats forts aux tests de Granger.

\paragraph{Fonds Propres - Non-Vie}\label{fonds-propres---non-vie}

\begin{equation}
\begin{aligned}
\Delta \ln(\text{FP}_t) = \alpha & + \rho \Delta \ln(\text{FP}_{t-1}) 
& + \beta_1 \Delta \ln(\text{IPC})_{t-3}
& + \beta_2 \Delta \text{TauxC}_{t-3}
& + \beta_3 \Delta \text{TauxL}_{t-3} + \epsilon_t
\end{aligned}
\end{equation}

\begin{table}[H]
\centering\centering
\caption{\label{tab:unnamed-chunk-25}Analyse dynamique des Fonds Propres - Branche Non-Vie}
\centering
\begin{tabular}[t]{lccc}
\toprule
\multicolumn{1}{c}{ } & \multicolumn{3}{c}{Fonds Propres (Non-Vie)} \\
\cmidrule(l{3pt}r{3pt}){2-4}
  & Est. & t & p\\
\midrule
\cellcolor{gray!10}{(Intercept)} & \cellcolor{gray!10}{\num{0.0367}**} & \cellcolor{gray!10}{\num{3.5278}} & \cellcolor{gray!10}{\num{0.0014}}\\
dplyr = lag(diff\_fonds\_propres, 1) & \num{-0.4123}*** & \num{-4.9547} & \num{<1e-04}\\
\cellcolor{gray!10}{dplyr = lag(diff\_ipc, 3)} & \cellcolor{gray!10}{\num{-4.1880}*} & \cellcolor{gray!10}{\num{-2.3698}} & \cellcolor{gray!10}{\num{0.0244}}\\
dplyr = lag(diff\_taux\_court, 3) & \num{-0.0332} & \num{-1.4900} & \num{0.1467}\\
\cellcolor{gray!10}{dplyr = lag(diff\_taux\_long, 3)} & \cellcolor{gray!10}{\num{0.0762}*} & \cellcolor{gray!10}{\num{2.1063}} & \cellcolor{gray!10}{\num{0.0436}}\\
\midrule
Num.Obs. & \num{35} &  & \\
\cellcolor{gray!10}{R2} & \cellcolor{gray!10}{\num{0.366}} & \cellcolor{gray!10}{} & \cellcolor{gray!10}{}\\
R2 Adj. & \num{0.281} &  & \\
\bottomrule
\end{tabular}
\end{table}

Les assureurs non-vie semblent davantage déconnectés des taux et des
marchés financiers que les assureurs Vie. La dynamique de leurs fonds
propres parait essentiellement dictée par sa propre inertie. Cependant,
la constante est significative indiquant une croissance structurelle
indépendante des fluctuations macro-financières. Cette dynamique suggère
que la richesse nette de cette branche repose principalement sur le
cycle d'exploitation (primes et sinistres) plutôt que sur le cycle
financier.

\paragraph{SCR - Non-Vie}\label{scr---non-vie}

\begin{equation}
\begin{aligned}
\Delta \ln(\text{SCR}_t) = \alpha & + \rho \Delta \ln(\text{SCR}_{t-1})
& + \beta_1 \Delta \text{TauxC}_{t-3} + \epsilon_t
\end{aligned}
\end{equation}

\begin{table}[H]
\centering\centering
\caption{\label{tab:unnamed-chunk-26}Analyse dynamique du SCR - Branche Non-Vie}
\centering
\begin{tabular}[t]{lccc}
\toprule
\multicolumn{1}{c}{ } & \multicolumn{3}{c}{SCR (Non-Vie)} \\
\cmidrule(l{3pt}r{3pt}){2-4}
  & Est. & t & p\\
\midrule
\cellcolor{gray!10}{(Intercept)} & \cellcolor{gray!10}{\num{0.0177}*} & \cellcolor{gray!10}{\num{2.2132}} & \cellcolor{gray!10}{\num{0.0341}}\\
dplyr = lag(diff\_scr, 1) & \num{-0.3273}*** & \num{-4.2606} & \num{0.0002}\\
\cellcolor{gray!10}{dplyr = lag(diff\_taux\_court, 3)} & \cellcolor{gray!10}{\num{-0.0233}**} & \cellcolor{gray!10}{\num{-3.0898}} & \cellcolor{gray!10}{\num{0.0041}}\\
\midrule
Num.Obs. & \num{35} &  & \\
\cellcolor{gray!10}{R2} & \cellcolor{gray!10}{\num{0.154}} & \cellcolor{gray!10}{} & \cellcolor{gray!10}{}\\
R2 Adj. & \num{0.101} &  & \\
\bottomrule
\end{tabular}
\end{table}

Les résultats montrent également qu'une hausse des taux d'intérêt de
court terme réduit la croissance du SCR avec un retard de trois
trimestres. Ce résultat suggère une transmission différée des conditions
de taux aux exigences réglementaires en capital, possiblement via les
mécanismes d'actualisation des engagements et l'évolution de
l'exposition au risque de taux. La constante positive suggère une
dynamique de croissance structurelle interne peu sensible aux chocs
macro-financiers immédiats.

\subsubsection{Ratio de solvabilité -
Non-Vie}\label{ratio-de-solvabilituxe9---non-vie}

\begin{equation}
\begin{aligned}
\Delta \ln(\text{Ratio\_Moyen}_t) = \alpha & + \rho \Delta \ln(\text{Ratio\_Moyen}_{t-1})
& + \beta_1 \Delta \text{TauxL}_{t-3} & + \beta_2 \Delta \text{TauxC}_{t-3} + \epsilon_t
\end{aligned}
\end{equation}

\begin{table}[H]
\centering\centering
\caption{\label{tab:unnamed-chunk-27}Analyse dynamique du ratio de solvabilité - Branche Non-Vie}
\centering
\begin{tabular}[t]{lccc}
\toprule
\multicolumn{1}{c}{ } & \multicolumn{3}{c}{Ratio (Non-Vie)} \\
\cmidrule(l{3pt}r{3pt}){2-4}
  & Est. & t & p\\
\midrule
\cellcolor{gray!10}{(Intercept)} & \cellcolor{gray!10}{\num{-0.0046}+} & \cellcolor{gray!10}{\num{-1.7371}} & \cellcolor{gray!10}{\num{0.0923}}\\
dplyr = lag(diff\_ratio, 1) & \num{-0.3611}*** & \num{-8.8205} & \num{<1e-04}\\
\cellcolor{gray!10}{dplyr = lag(diff\_taux\_court, 3)} & \cellcolor{gray!10}{\num{-0.0273}***} & \cellcolor{gray!10}{\num{-6.4019}} & \cellcolor{gray!10}{\num{<1e-04}}\\
dplyr = lag(diff\_taux\_long, 3) & \num{0.0407}** & \num{3.1076} & \num{0.0040}\\
\midrule
\cellcolor{gray!10}{Num.Obs.} & \cellcolor{gray!10}{\num{35}} & \cellcolor{gray!10}{} & \cellcolor{gray!10}{}\\
R2 & \num{0.253} &  & \\
\cellcolor{gray!10}{R2 Adj.} & \cellcolor{gray!10}{\num{0.180}} & \cellcolor{gray!10}{} & \cellcolor{gray!10}{}\\
\bottomrule
\end{tabular}
\end{table}

Une hausse des taux d'intérêt de long terme améliore le ratio de
solvabilité avec un retard de trois trimestres, probablement via
l'augmentation des fonds propres et l'amélioration des revenus
financiers. À l'inverse, les taux courts exercent un effet négatif sur
le ratio de solvabilité. Ce résultat suggère qu'un relèvement des taux
courts peut temporairement détériorer l'équilibre prudentiel des
assureurs Non-Vie.

\subsubsection{5) Diagnostics et limites du
modèle}\label{diagnostics-et-limites-du-moduxe8le}

Ces résultats doivent être interprétés avec prudence en raison de
plusieurs limites économétriques. La taille réduite de l'échantillon
restreint la puissance statistique des estimations. L'ensemble des
modèles estimés satisfont les conditions de validité économétrique : les
tests de Jarque-Bera, Breusch-Pagan et Ljung-Box ne rejettent pas la
normalité des résidus, l'homoscédasticité et l'absence d'autocorrélation
(cf.~Annexe Tables). Néanmoins, compte tenu de la faible puissance de
ces tests sur petit échantillon et de la sensibilité des p-values à la
correction, l'estimateur de Newey-West a été appliqué à l'ensemble des
modèles afin de garantir la robustesse des inférences. Par ailleurs, un
biais d'endogénéité lié à l'omission de variables pertinentes ne peut
être exclu. Les variables macro-financières retenues sont susceptibles
d'être corrélées avec d'autres déterminants de la solvabilité non
intégrés au modèle ce qui limite la validité de cette modélisation.

Par ailleurs, cette approche en différences ne permet pas d'identifier
les relations d'équilibre de long terme. Ces éléments justifient le
recours à une spécification ARDL qui permet d'intégrer simultanément
dynamiques de court terme et relations de cointégration éventuelles.

\section{Conclusion}\label{conclusion}

\section{Annexes}\label{annexes}

\subsubsection{Analyse descriptive}\label{analyse-descriptive}

\begin{flushleft}\includegraphics{Stat_desc_files/figure-latex/unnamed-chunk-28-1} \end{flushleft}

\begin{flushleft}\includegraphics{Stat_desc_files/figure-latex/unnamed-chunk-30-1} \end{flushleft}

\subsubsection{Les indices boursiers.}\label{les-indices-boursiers.}

Les indices boursiers reflètent la performance globale des actions. Une
baisse de ces indices entraîne une diminution de la valeur des actifs en
actions détenus par les assureurs, ce qui se traduit par une
augmentation du SCR actions et une réduction du ratio de solvabilité.
Dans cette étude, deux indices boursiers sont considérés: l'EURO STOXX
et le CAC 40.

\paragraph{L'EURO STOXX 50}\label{leuro-stoxx-50}

L'\emph{EURO STOXX 50} est un indice boursier regroupant les 50
principales sociétés cotées de la zone euro, sélectionnées
principalement en fonction de leur capitalisation boursière flottante.
Il constitue un indicateur de la santé économique et des performances
des marchés financiers au sein de la zone euro. Le graphique présenté
ci-dessous illustre l'évolution des valeurs maximales et minimales de
cet indice deouis \(2015\) jusqu'en \(2025\).

\begin{flushleft}\includegraphics{Stat_desc_files/figure-latex/unnamed-chunk-34-1} \end{flushleft}

Ces indices évoluent de manière globalement parallèle, suivant une
tendance générale à la hausse. Néanmoins, une chute marquée est observée
en mars 2020, atteignant un niveau minimal, en raison des mesures de
restriction et de l'incertitude économique engendrées par la pandémie de
COVID-19. C'est en effet en cette période que Donald Trump avait annoncé
la restriction des entrées sur le téritoire américain; source d'une
panique boursière.

\paragraph{Le Standard \& Poor's 500 (S\&P
500)}\label{le-standard-poors-500-sp-500}

Le S\&P 500 est un indice boursier large représentant les 500 plus
grandes sociétés cotées aux États-Unis, pondérées en fonction de leur
capitalisation boursière flottante. Il est largement utilisé comme
référence du marché actions américain et constitue un indicateur clé de
la performance économique et financière du pays. L'évolution des valeurs
maximales et minimales de cet indice entre 2015 et 2025, présentée dans
le graphique ci-dessous, permet d'analyser la volatilité du marché ainsi
que l'impact des événements macroéconomiques et géopolitiques sur les
marchés financiers américains

\begin{flushleft}\includegraphics{Stat_desc_files/figure-latex/unnamed-chunk-35-1} \end{flushleft}

Le S\&P 500 connaît une tendence croissance sur toute la période avec
également une pertubation en mars 2020 due à la pandémie de OVID-19.
Toutefois cette indice est visiblement moins volatile que l'EURO STOXX
50.

\paragraph{L'indice CAC 40}\label{lindice-cac-40}

~

L'indice \emph{CAC 40} constitue le principal indice boursier de la
Bourse de Paris, reflétant la performance des 40 actions les plus
importantes et les plus activement négociées sur Euronext Paris. Le
graphique ci-après illustre l'évolution de ses valeurs maximales et
minimales au fil du temps.

\begin{flushleft}\includegraphics{Stat_desc_files/figure-latex/unnamed-chunk-38-1} \end{flushleft}

Comme pour l'EURO STOXX, les valeurs de l'indice CAC 40 évoluent de
manière parallèle et croissante avec également une chutte marqué en mars
2020 pour la même raison que précédemment.

\paragraph{L'indice DAX}\label{lindice-dax}

Le DAX (Deutscher Aktienindex) est l'indice de référence du marché
boursier allemand, regroupant les 40 principales sociétés cotées à la
Bourse de Francfort, sélectionnées principalement en fonction de leur
capitalisation boursière flottante. C'est l'équivalent allemend du CAC
40. Le graphique ci-après représente les variation de cet indice sur la
période \(2015-2025.\)

\begin{flushleft}\includegraphics{Stat_desc_files/figure-latex/unnamed-chunk-39-1} \end{flushleft}

\section{Modélisation}\label{moduxe9lisation-1}

\section{II - Régression en
différence}\label{ii---ruxe9gression-en-diffuxe9rence}

\begin{table}[H]
\centering\centering
\caption{\label{tab:unnamed-chunk-40}Résultats de la régression du ratio de solvabilité - Vie}
\centering
\begin{tabular}[t]{lcccc}
\toprule
\multicolumn{1}{c}{ } & \multicolumn{4}{c}{(1)} \\
\cmidrule(l{3pt}r{3pt}){2-5}
  & Est. & S.E. & t & p\\
\midrule
\cellcolor{gray!10}{(Intercept)} & \cellcolor{gray!10}{\num{-0.0158}} & \cellcolor{gray!10}{\num{0.0205}} & \cellcolor{gray!10}{\num{-0.7704}} & \cellcolor{gray!10}{\num{0.4475}}\\
rendement\_stoxx & \num{0.1268} & \num{0.2570} & \num{0.4936} & \num{0.6255}\\
\cellcolor{gray!10}{rendement\_sp500} & \cellcolor{gray!10}{\num{0.2066}} & \cellcolor{gray!10}{\num{0.2412}} & \cellcolor{gray!10}{\num{0.8566}} & \cellcolor{gray!10}{\num{0.3989}}\\
diff\_ipc & \num{0.6865} & \num{2.6065} & \num{0.2634} & \num{0.7942}\\
\cellcolor{gray!10}{diff\_taux\_court} & \cellcolor{gray!10}{\num{-0.0432}} & \cellcolor{gray!10}{\num{0.0382}} & \cellcolor{gray!10}{\num{-1.1325}} & \cellcolor{gray!10}{\num{0.2670}}\\
diff\_taux\_long & \num{0.0599} & \num{0.0455} & \num{1.3172} & \num{0.1984}\\
\cellcolor{gray!10}{diff\_pib} & \cellcolor{gray!10}{\num{0.2577}} & \cellcolor{gray!10}{\num{0.4199}} & \cellcolor{gray!10}{\num{0.6138}} & \cellcolor{gray!10}{\num{0.5443}}\\
\midrule
Num.Obs. & \num{35} &  &  & \\
\cellcolor{gray!10}{R2} & \cellcolor{gray!10}{\num{0.176}} & \cellcolor{gray!10}{} & \cellcolor{gray!10}{} & \cellcolor{gray!10}{}\\
R2 Adj. & \num{-0.000} &  &  & \\
\bottomrule
\end{tabular}
\end{table}

\begin{longtable}[t]{llrrl}
\caption{\label{tab:unnamed-chunk-41}Tests sur les résidus modèle ratio - Vie}\\
\toprule
 & Test & Statistique & P\_value & Verdict\\
\midrule
\cellcolor{gray!10}{DW} & \cellcolor{gray!10}{Autocorrélation (Durbin-Watson)} & \cellcolor{gray!10}{2.6972} & \cellcolor{gray!10}{0.9789} & \cellcolor{gray!10}{Hypothèse validée}\\
X-squared & Indépendance Globale (Ljung-Box) & 17.6091 & 0.0619 & Bruit blanc confirmé\\
\cellcolor{gray!10}{BP} & \cellcolor{gray!10}{Homoscédasticité (Breusch-Pagan)} & \cellcolor{gray!10}{4.5900} & \cellcolor{gray!10}{0.5974} & \cellcolor{gray!10}{Hypothèse validée}\\
W & Normalité (Shapiro-Wilk) & 0.9649 & 0.3184 & Hypothèse validée\\
\bottomrule
\end{longtable}

\begin{longtable}[t]{lrrl}
\caption{\label{tab:unnamed-chunk-42}Tests diagnostiques - Fonds Propres Vie}\\
\toprule
Test & Statistique & p.value & Hypothèse\_nulle\\
\midrule
\cellcolor{gray!10}{Ljung-Box} & \cellcolor{gray!10}{0.0971} & \cellcolor{gray!10}{0.7554} & \cellcolor{gray!10}{Pas d'autocorrélation}\\
Jarque-Bera & 0.1760 & 0.9158 & Normalité des résidus\\
\cellcolor{gray!10}{Breusch-Pagan} & \cellcolor{gray!10}{1.7255} & \cellcolor{gray!10}{0.7861} & \cellcolor{gray!10}{Homoscédasticité}\\
\bottomrule
\end{longtable}

\begin{longtable}[t]{llr}
\caption{\label{tab:unnamed-chunk-42}Facteurs d'inflation de la variance (VIF) - Fonds Propres Vie}\\
\toprule
 & Variable & VIF\\
\midrule
\cellcolor{gray!10}{dplyr::lag(diff\_fonds\_propres, 1)} & \cellcolor{gray!10}{dplyr::lag(diff\_fonds\_propres, 1)} & \cellcolor{gray!10}{1.309}\\
dplyr::lag(diff\_taux\_long, 1) & dplyr::lag(diff\_taux\_long, 1) & 1.051\\
\cellcolor{gray!10}{dplyr::lag(rendement\_sp500, 2)} & \cellcolor{gray!10}{dplyr::lag(rendement\_sp500, 2)} & \cellcolor{gray!10}{1.068}\\
dplyr::lag(diff\_ipc, 4) & dplyr::lag(diff\_ipc, 4) & 1.281\\
\bottomrule
\end{longtable}

\begin{longtable}[t]{lrrl}
\caption{\label{tab:unnamed-chunk-43}Tests diagnostiques - SCR Vie}\\
\toprule
Test & Statistique & p.value & Hypothèse\_nulle\\
\midrule
\cellcolor{gray!10}{Ljung-Box} & \cellcolor{gray!10}{0.0971} & \cellcolor{gray!10}{0.7554} & \cellcolor{gray!10}{Pas d'autocorrélation}\\
Jarque-Bera & 0.1760 & 0.9158 & Normalité des résidus\\
\cellcolor{gray!10}{Breusch-Pagan} & \cellcolor{gray!10}{1.7255} & \cellcolor{gray!10}{0.7861} & \cellcolor{gray!10}{Homoscédasticité}\\
\bottomrule
\end{longtable}

\begin{longtable}[t]{llr}
\caption{\label{tab:unnamed-chunk-43}Facteurs d'inflation de la variance (VIF) - SCR Vie}\\
\toprule
 & Variable & VIF\\
\midrule
\cellcolor{gray!10}{dplyr::lag(diff\_scr, 1)} & \cellcolor{gray!10}{dplyr::lag(diff\_scr, 1)} & \cellcolor{gray!10}{1.054}\\
dplyr::lag(diff\_taux\_long, 1) & dplyr::lag(diff\_taux\_long, 1) & 1.212\\
\cellcolor{gray!10}{dplyr::lag(diff\_ipc, 3)} & \cellcolor{gray!10}{dplyr::lag(diff\_ipc, 3)} & \cellcolor{gray!10}{1.243}\\
\bottomrule
\end{longtable}

\begin{longtable}[t]{lrrl}
\caption{\label{tab:unnamed-chunk-44}Tests diagnostiques - Ratio de solvabilité Vie}\\
\toprule
Test & Statistique & p.value & Hypothèse\_nulle\\
\midrule
\cellcolor{gray!10}{Ljung-Box} & \cellcolor{gray!10}{0.6204} & \cellcolor{gray!10}{0.4309} & \cellcolor{gray!10}{Pas d'autocorrélation}\\
Jarque-Bera & 4.2546 & 0.1192 & Normalité des résidus\\
\cellcolor{gray!10}{Breusch-Pagan} & \cellcolor{gray!10}{3.5422} & \cellcolor{gray!10}{0.1701} & \cellcolor{gray!10}{Homoscédasticité}\\
\bottomrule
\end{longtable}

\begin{longtable}[t]{llr}
\caption{\label{tab:unnamed-chunk-44}Facteurs d'inflation de la variance (VIF) - Ratio Vie}\\
\toprule
 & Variable & VIF\\
\midrule
\cellcolor{gray!10}{dplyr::lag(diff\_ratio, 1)} & \cellcolor{gray!10}{dplyr::lag(diff\_ratio, 1)} & \cellcolor{gray!10}{1}\\
dplyr::lag(rendement\_sp500, 3) & dplyr::lag(rendement\_sp500, 3) & 1\\
\bottomrule
\end{longtable}

\begin{longtable}[t]{lrrl}
\caption{\label{tab:unnamed-chunk-45}Tests diagnostiques - Fonds propres Non-Vie}\\
\toprule
Test & Statistique & p.value & Hypothèse\_nulle\\
\midrule
\cellcolor{gray!10}{Ljung-Box} & \cellcolor{gray!10}{0.3490} & \cellcolor{gray!10}{0.5547} & \cellcolor{gray!10}{Pas d'autocorrélation}\\
Jarque-Bera & 3.7717 & 0.1517 & Normalité des résidus\\
\cellcolor{gray!10}{Breusch-Pagan} & \cellcolor{gray!10}{1.3123} & \cellcolor{gray!10}{0.8593} & \cellcolor{gray!10}{Homoscédasticité}\\
\bottomrule
\end{longtable}

\begin{longtable}[t]{llr}
\caption{\label{tab:unnamed-chunk-45}Facteurs d'inflation de la variance (VIF) - Fond propres Non-Vie}\\
\toprule
 & Variable & VIF\\
\midrule
\cellcolor{gray!10}{dplyr::lag(diff\_fonds\_propres, 1)} & \cellcolor{gray!10}{dplyr::lag(diff\_fonds\_propres, 1)} & \cellcolor{gray!10}{1.055}\\
dplyr::lag(diff\_ipc, 3) & dplyr::lag(diff\_ipc, 3) & 1.465\\
\cellcolor{gray!10}{dplyr::lag(diff\_taux\_court, 3)} & \cellcolor{gray!10}{dplyr::lag(diff\_taux\_court, 3)} & \cellcolor{gray!10}{1.622}\\
dplyr::lag(diff\_taux\_long, 3) & dplyr::lag(diff\_taux\_long, 3) & 1.325\\
\bottomrule
\end{longtable}

\begin{longtable}[t]{lrrl}
\caption{\label{tab:unnamed-chunk-46}Tests diagnostiques - Fonds propres Non-Vie}\\
\toprule
Test & Statistique & p.value & Hypothèse\_nulle\\
\midrule
\cellcolor{gray!10}{Ljung-Box} & \cellcolor{gray!10}{0.0381} & \cellcolor{gray!10}{0.8452} & \cellcolor{gray!10}{Pas d'autocorrélation}\\
Jarque-Bera & 18.0553 & 0.0001 & Normalité des résidus\\
\cellcolor{gray!10}{Breusch-Pagan} & \cellcolor{gray!10}{0.2348} & \cellcolor{gray!10}{0.8892} & \cellcolor{gray!10}{Homoscédasticité}\\
\bottomrule
\end{longtable}

\begin{longtable}[t]{llr}
\caption{\label{tab:unnamed-chunk-46}Facteurs d'inflation de la variance (VIF) - Fond propres Non-Vie}\\
\toprule
 & Variable & VIF\\
\midrule
\cellcolor{gray!10}{dplyr::lag(diff\_scr, 1)} & \cellcolor{gray!10}{dplyr::lag(diff\_scr, 1)} & \cellcolor{gray!10}{1.012}\\
dplyr::lag(diff\_taux\_court, 3) & dplyr::lag(diff\_taux\_court, 3) & 1.012\\
\bottomrule
\end{longtable}

\begin{longtable}[t]{lrrl}
\caption{\label{tab:unnamed-chunk-47}Tests diagnostiques - Fonds propres Non-Vie}\\
\toprule
Test & Statistique & p.value & Hypothèse\_nulle\\
\midrule
\cellcolor{gray!10}{Ljung-Box} & \cellcolor{gray!10}{0.0211} & \cellcolor{gray!10}{0.8846} & \cellcolor{gray!10}{Pas d'autocorrélation}\\
Jarque-Bera & 101.6076 & 0.0000 & Normalité des résidus\\
\cellcolor{gray!10}{Breusch-Pagan} & \cellcolor{gray!10}{2.7860} & \cellcolor{gray!10}{0.4258} & \cellcolor{gray!10}{Homoscédasticité}\\
\bottomrule
\end{longtable}

\begin{longtable}[t]{llr}
\caption{\label{tab:unnamed-chunk-47}Facteurs d'inflation de la variance (VIF) - Fond propres Non-Vie}\\
\toprule
 & Variable & VIF\\
\midrule
\cellcolor{gray!10}{dplyr::lag(diff\_ratio, 1)} & \cellcolor{gray!10}{dplyr::lag(diff\_ratio, 1)} & \cellcolor{gray!10}{1.083}\\
dplyr::lag(diff\_taux\_court, 3) & dplyr::lag(diff\_taux\_court, 3) & 1.314\\
\cellcolor{gray!10}{dplyr::lag(diff\_taux\_long, 3)} & \cellcolor{gray!10}{dplyr::lag(diff\_taux\_long, 3)} & \cellcolor{gray!10}{1.321}\\
\bottomrule
\end{longtable}

\end{document}
